\documentclass[11pt,a4paper]{article}

\usepackage[margin=24mm]{geometry}
\usepackage[T1]{fontenc}
\usepackage[utf8]{inputenc}
\usepackage{lmodern}
\usepackage{microtype}
\usepackage{booktabs}
\usepackage{longtable}
\usepackage{tabularx}
\usepackage{array}
\usepackage{enumitem}
\usepackage{xcolor}
\usepackage{listings}
\usepackage{hyperref}
\usepackage{fancyhdr}

\hypersetup{
  colorlinks=true,
  linkcolor=blue!40!black,
  urlcolor=blue!50!black,
  citecolor=blue!40!black,
  pdftitle={HEPLab Project Documentation},
  pdfauthor={HEPLab Project}
}
\urlstyle{tt}
\setlist{itemsep=2pt,topsep=4pt}
\setcounter{tocdepth}{2}
\newcolumntype{Y}{>{\raggedright\arraybackslash}X}
\newcolumntype{P}[1]{>{\raggedright\arraybackslash}p{#1}}

\lstset{
  basicstyle=\ttfamily\small,
  columns=fullflexible,
  breaklines=true,
  frame=single,
  rulecolor=\color{black!20},
  backgroundcolor=\color{black!3},
  showstringspaces=false,
  keepspaces=true
}

\pagestyle{fancy}
\fancyhf{}
\lhead{HEPLab}
\rhead{Project documentation}
\cfoot{\thepage}
\setlength{\headheight}{14pt}

\title{\textbf{HEPLab}\\[4pt]\large Project Documentation}
\author{Project overview, architecture, scientific scope, and current status}
\date{23 November 2026}

\begin{document}
\maketitle

\begin{quote}
\textbf{Status and scope.} This document summarizes the HEPLab repository as
reviewed on 23 November 2026. It is a guide to the project and its current
evidence, not a claim that every planned component is complete or that the
whole application is a validated detector simulation. For detailed algorithms
and release gates, use the engineering handbook cited at the end.
\end{quote}

\tableofcontents
\clearpage

\section{Executive overview}

HEPLab is an open-source educational high-energy physics (HEP) workbench. It
brings together selected workflows for reading CMS Open Data, performing
dimuon analyses, running a fast Monte Carlo simulation chain, exploring
relativistic particle motion, and teaching analysis concepts through notebooks
and an offline Z-boson lesson.

The project has three scientific paths that must be kept distinct:

\begin{enumerate}
  \item \textbf{Recorded CMS data:} real NanoAOD inputs and analyses with
  source, certification, selection, and fit provenance.
  \item \textbf{External fast simulation:} generated events passed through a
  configured PYTHIA8 and DELPHES chain. Its predictions depend on the exact
  process and detector card.
  \item \textbf{Analytic educational models:} controlled calculations such as
  four-vector operations, ideal-field trajectories, and simplified material
  response.
\end{enumerate}

These paths can support useful teaching and software comparisons, but they do
not make HEPLab a globally validated CMS detector model. A plot, a successful
software check, or a comparison with one reference sample establishes only the
specific behavior and domain checked.

\subsection{Project purpose}

HEPLab aims to make particle-physics workflows easier to inspect and reproduce.
A learner or analyst should be able to follow a result from its input dataset
or simulation settings through object selection and statistical summaries to
the saved output and its provenance. The platform is intended for students,
educators, analysts, and open-source contributors.

\subsection{Current maturity}

The repository describes the product as an early implementation. It contains a
Python analysis and workbench core, a parity-checked Rust numerical backend, a
PYTHIA8/DELPHES integration path, educational physics models, a web guide, and
an evolving native desktop application. Important validation, calibration,
integration, accessibility, and release tasks remain open. In particular,
the project does not claim that any requirement has reached a general
\emph{physics-validated} state.

\section{Goals and boundaries}

\subsection{Goals}

The project is organized around the following goals:

\begin{itemize}
  \item Provide a practical route from a dataset catalog and event input to
  analysis outputs such as cutflows, histograms, fits, and plots.
  \item Give users a clear way to distinguish measured input, generated
  simulation, educational fixtures, model predictions, and validation
  evidence.
  \item Keep data and simulation workflows inspectable through versioned
  manifests, explicit event contracts, reproducible configuration, and
  preserved outputs.
  \item Teach core ideas in HEP analysis and detector response without
  presenting simplified models as complete experimental software.
  \item Support a local-first workflow that can be developed and demonstrated
  without a hosted multi-user service.
\end{itemize}

\subsection{Boundaries}

The following are outside the project's current validated scope:

\begin{itemize}
  \item A full CMS reconstruction chain or a faithful simulation of every CMS
  subsystem and data-taking condition.
  \item A general-purpose Geant4 transport replacement, a complete atmospheric
  muon shower simulation, or detector response calibrated across all species,
  energies, and geometries.
  \item Quantum-state, spin, entanglement, or weak-decay-amplitude simulation.
  \item A new precision measurement or a claim of new physics.
  \item A production hosted service, signed cross-platform desktop release, or
  complete native desktop integration.
  \item General support for experiments other than CMS.
\end{itemize}

\section{Product workflows}

\subsection{CMS Open Data analysis}

The data workflow starts from a named dataset and a frozen input selection.
Readers normalize NanoAOD records into the project's event representation;
the analysis applies its declared object and event selections; and outputs
include analysis products and provenance. The repository includes a catalog
entry for the 2016G DoubleMuon sample and a verified 100-event remote-read
demonstration. Its NanoAOD path supports chunked reads, source checks, cached
data, and certification support.

The principal analysis path is the dimuon spectrum. Opposite-charge muons are
selected according to an explicit policy, their invariant mass is calculated,
and the workflow produces cutflows, histogram data, selected-pair information,
and fit or plot outputs as supported by the selected command. A run must be
interpreted in the context of its input coverage, trigger policy, cuts,
weights, corrections, and fit model.

\subsection{Generation and fast detector simulation}

A separate simulation path uses PYTHIA8 for event generation and DELPHES for
fast detector simulation. The integration is available through a registered
runtime adapter and has been smoke-tested in Linux/WSL environments. A
generator process, tune, PDF, cuts, detector card, normalization, and software
versions all affect its output.

This is a fast-simulation reference path. It is not the same as full CMS
detector transport and does not by itself validate agreement with recorded
data. General process coverage, pileup, a complete correction inventory,
matched reference samples, and independent cross-section and detector-response
validation remain delivery work.

\subsection{Analytic educational workbench}

The educational physics path provides selected, bounded calculations. These
include relativistic four-vectors, ideal-solenoid motion, selected geometry
crossings, and simplified material or electronics response. Users can inspect
predicted trajectories and model outputs as teaching aids.

The displayed crossings are model predictions, not measured detector hits.
The material examples do not constitute complete material transport. A
separate fixed-momentum vacuum muon-survival exercise is not coupled to the
track simulation and does not model atmospheric production or shower
development.

\subsection{Learning and event visualization}

The project includes an offline Z-boson learning workflow and a gallery of
notebooks. Exercises guide users through event inspection and analysis while
identifying synthetic fixtures and the limits of each result. Training
materials and replay-oriented laboratory modules extend the educational
workflow.

Event displays are intended to help users inspect a selected event and its
model context. The project can export event views as vendor-neutral OpenUSD
stages for offline use in compatible viewers. An event display is a
visualization, not a substitute for the analysis artifact, selection record,
or scientific validation.

\section{Architecture}

\subsection{System outline}

The repository is a multi-component project. Python provides most current
analysis and application services; Rust provides selected numerical kernels;
C++ and Qt are being used for the native desktop rebuild; and external
generator and detector packages provide a separate fast-simulation route.

\begin{center}
\small
\setlength{\tabcolsep}{5pt}
\begin{tabularx}{\textwidth}{P{0.24\textwidth} Y Y}
\toprule
\textbf{Component} & \textbf{Role} & \textbf{Current boundary} \\
\midrule
Python package &
Event readers, analysis, physics reference code, CLI, local service, and
workbench operations. &
The main integration layer; optional analysis dependencies are installed
separately. \\
Rust numerical core &
Selected physics kernels, exposed to Python through PyO3. &
Reference and parity backend for a bounded scope; no general speedup claim. \\
PYTHIA8 and DELPHES &
External event generation and fast detector simulation. &
Card- and process-specific; overall physics validation is incomplete. \\
Native desktop &
C++20 and Qt 6 Workbench rebuild with a native application shell. &
Work in progress; the shell is not yet fully integrated with all scientific
services or release targets. \\
ASP.NET Core site &
Public project guide and educational website, using Razor Pages and .NET 10. &
Separate from the local scientific workbench. \\
FastAPI local surface &
Compatibility GUI and local service operations for selected workflows. &
Not the native Qt application and not a hosted collaboration service. \\
Jupyter and MkDocs &
Notebook learning material and the engineering handbook. &
Documentation and educational assets; generated site output is not the
canonical source. \\
\bottomrule
\end{tabularx}
\end{center}

\subsection{Repository map}

\begin{longtable}{P{0.28\textwidth} P{0.64\textwidth}}
\toprule
\textbf{Path} & \textbf{Contents} \\
\midrule
\endhead
\path{src/heplab/} & Python package: physics, readers, analyses, simulation,
I/O, schemas, server, laboratory replay, and exports. \\
\path{kernels/} & Rust numerical core and Python bindings. \\
\path{desktop/} & Native desktop application sources and build configuration. \\
\path{web/HEPLab.Web/} & ASP.NET Core project guide and learning website. \\
\path{notebooks/} & Guided notebooks, including the Z-boson masterclass and
gallery. \\
\path{docs/engineering/} & Engineering handbook, requirements, implementation
evidence, and delivery plans. \\
\path{docs/scientific-review/} & Scientific scope, verification, review,
and presentation guidance. \\
\path{contracts/} & JSON schemas and example data contracts. \\
\path{manifests/} & Dataset and laboratory input manifests. \\
\path{tests/} & Python and integration-level regression tests and fixtures. \\
\path{tools/} & Data acquisition, build helpers, verification tools, and
documentation checks. \\
\path{artifacts/} & Preserved example and run outputs; artifacts should be
interpreted with their accompanying provenance. \\
\bottomrule
\end{longtable}

\subsection{Typical data flow}

\begin{center}
\fbox{\parbox{0.92\textwidth}{\centering
Dataset manifest or simulation configuration
$\rightarrow$ verified input and provenance
$\rightarrow$ reader and event contract
$\rightarrow$ selection and analysis
$\rightarrow$ histogram, fit, or event view
$\rightarrow$ saved result and evidence status}}
\end{center}

The same high-level analysis concepts can be applied to different event
sources, but shared field names do not guarantee that the source objects have
equivalent meanings. For example, an isolation value or trigger decision must
retain its definition and availability. Incompatible inputs should be rejected
or handled through an explicitly reviewed mapping.

\section{Data and analysis methods}

\subsection{Inputs and contracts}

The current project uses schema-v1 event normalization in its established
reader paths. Contracts under \path{contracts/} define additional and
proposed interfaces, including dataset, detector, run-plan, and event-chunk
structures. The proposed event schema v2 is a target for a versioned
migration; it should not be described as the schema used by every existing
artifact.

The CMS data path uses dataset manifests, file identities, source entry ranges,
and certified-luminosity information where configured. NanoAOD is read in
chunks so that processing need not hold an entire large dataset in memory.
Cached event content and the analysis configuration should remain associated
with the exact input and selection used.

\subsection{Analysis and statistics}

The dimuon workflow covers candidate selection, mass calculation, weighted
histograms or cutflows where configured, and fit operations using the available
CPU statistics stack, including SciPy and iminuit. The interpretation of a fit
depends on its range, background and signal model, normalization, candidate
policy, and uncertainty treatment.

A simulated sample's normalization is tied to its cross section, integrated
luminosity, filter efficiency, and generator-weight convention. Negative
generator weights must be preserved where present. Data and simulation should
not be compared as though they were interchangeable: era, object definitions,
triggers, weights, and correction compatibility all matter.

The data-analysis workflow is useful for reproducible exploration and
teaching. Precision efficiency, trigger, momentum-scale, and systematic
uncertainties require further sourced inputs, matched samples, and review
before supporting measurement claims.

\subsection{Provenance and saved results}

A useful result should identify its input dataset or simulation settings, event
range or sample scope, relevant configuration, software/model versions,
selection and fit choices, and validation status. The project's engineering
direction favors content-addressed inputs and immutable result artifacts.
Current workflows provide provenance features, but the engineering handbook
records further work needed to capture a complete source snapshot, all
correction payloads, random state, runtime details, and output identity in a
fully preserved bundle.

\section{Scientific models and evidence}

\subsection{Model domains}

\begin{longtable}{P{0.24\textwidth} P{0.36\textwidth} P{0.32\textwidth}}
\toprule
\textbf{Model or workflow} & \textbf{Evidence available} & \textbf{Limit of the claim} \\
\midrule
\endhead
Four-vector and kinematics &
Analytic reference calculations with selected independent checks and Rust
parity evidence. &
Applies to the checked primitives and their stated numeric domains; it does not
validate a detector model. \\
Uniform-field motion &
Closed-form ideal-solenoid trajectories and comparisons for selected
field-translation calculations. &
Idealized or bounded field models; crossings are predictions, not measured
tracks. \\
Geometry and surfaces &
Selected cylinders, disks, and sourced display geometry are implemented. &
Not complete CMS geometry, alignment, conditions, or material navigation. \\
Silicon energy loss &
A sourced mean-stopping table fit has a held-out table check. &
A mean response check is not full stochastic transport or detector
calibration. \\
Material and electronics &
Educational thin-slab response and deposit-to-ADC examples. &
Simplified validity domain; full stochastic interactions, pileup, and calibrated
electronics response remain open. \\
Muon response transfer &
A coarse DELPHES-based transfer model has bounded reference checks. &
Not an experimental CMS calibration or a universally valid muon response. \\
PYTHIA8--DELPHES &
A runtime adapter and smoke-tested fast-simulation chain. &
Results depend on the process and card; this is not full CMS simulation or
general closure evidence. \\
CMS resonance regression &
A mass-independent candidate policy and event-disjoint split are used in the
current low-mass regression. &
The split is from one source file, and the frozen aggregate gate fails at
$\Upsilon(2S)$. It is not independent-dataset validation. \\
Muon survival lesson &
An isolated fixed-momentum vacuum survival calculation is available. &
Not an atmospheric muon, air-shower, or coupled detector-transport model. \\
\bottomrule
\end{longtable}

\subsection{Fidelity language}

The engineering handbook distinguishes analytic/toy calculations, calibrated
parametric detector response, DELPHES fast simulation, and optional detailed
transport such as Geant4. These levels answer different questions. A model may
be helpful at one level and still lack evidence at another. In particular:

\begin{itemize}
  \item An analytic formula can be verified against a closed form without
  establishing that its assumptions describe the CMS detector.
  \item A software-reference comparison can show agreement with a pinned
  implementation without establishing agreement with experimental data.
  \item A data regression can check a named sample and selection without
  validating a general simulation or a full reconstruction chain.
  \item Passing software tests establishes behavior covered by those tests; it
  does not prove physical accuracy outside the checked cases.
\end{itemize}

\subsection{Current scientific assessment}

The defensible description is an educational simulator with specific checked
submodels and an early-stage HEP data-analysis workflow. Some kinematic,
ideal-field, material-table, and field-reference calculations have bounded
evidence. The analytic detector is not a complete transport or reconstruction
model.

The current low-mass CMS regression corrects a prior target-dependent
candidate-selection problem and splits events before fitting. Its aggregate
reference gate still fails at the $\Upsilon(2S)$ peak, and its training and
held-out events come from the same source file. The data result is therefore
exploratory and must not be described as independent experimental calibration.

\section{Verification and release practice}

HEPLab maintains a traceability approach connecting requirements, contracts,
implementations, acceptance cases, and evidence artifacts. Verification may
include structural contract checks, analytic numerical references, statistical
checks, reader fixtures, integration cases, fault recovery, packaging, and
physics comparisons. These evidence types are complementary rather than
interchangeable.

A release claim should identify:
\begin{itemize}
  \item the exact model, software version, configuration, and valid domain;
  \item the data or simulation identity and event coverage;
  \item the independent reference or acceptance case used;
  \item the result, tolerance, and uncertainty treatment; and
  \item unresolved limitations or failed gates that affect interpretation.
\end{itemize}

Documentation and schema checkers can detect broken links, malformed
contracts, or inconsistent identifiers. They cannot establish that an
algorithm is scientifically correct. Likewise, a successful build or test suite
does not amount to independent physics review.

\section{Current implementation status}

\begin{longtable}{P{0.24\textwidth} P{0.38\textwidth} P{0.30\textwidth}}
\toprule
\textbf{Area} & \textbf{Current state} & \textbf{Open work or boundary} \\
\midrule
\endhead
CMS data &
NanoAOD reader, chunked data path, certification support, range-aware cache,
and a 100-event remote-read demonstration. &
Production-scale throughput, broad independent samples, full correction
inventory, and complete preservation evidence. \\
Dimuon analysis &
Candidate selection, mass spectra, cutflows, fits, and an exploratory CMS
reference workflow. &
The low-mass aggregate gate fails at $\Upsilon(2S)$; trigger, efficiency,
correction, and systematic evidence remain limited. \\
Rust core &
Numerical core and PyO3 binding with parity checks for bounded primitives. &
No general performance claim; broader platform packaging remains a release
gate. \\
External simulation &
PYTHIA8--DELPHES integration and runtime smoke path. &
General processes, matched reference MC, pileup, and independent physics
validation. \\
Educational physics &
Analytic motion, selected geometry and material models, survival lesson,
event inspection, and offline Z-boson learning material. &
The model domains are intentionally limited and must remain visible to users. \\
Native desktop &
C++20/Qt 6 native Workbench rebuild is underway. &
Scientific-service integration, persistence, platform acceptance, installer,
and signing work remain. \\
Website and handbook &
ASP.NET Core project guide and MkDocs-authored engineering handbook. &
Public deployment and release ownership are separate from local documentation
builds. \\
Laboratory replay &
Device-profile, event, acquisition, and replay-oriented project components. &
A supported physical device and its end-to-end measurement validation need
their own evidence. \\
\bottomrule
\end{longtable}

\section{Priorities and roadmap}

The project should progress according to evidence and dependencies, rather than
feature count alone. The near-term priorities documented in the repository are:

\begin{enumerate}
  \item Strengthen the CMS data foundation with appropriate independent
  samples, correction and trigger provenance, and exact source coverage.
  \item Resolve and report the failing low-mass resonance regression using
  predeclared, mass-independent selection and independent data where available.
  \item Match CMS data with a suitable official simulation sample before making
  detector-response or closure claims.
  \item Continue calibration and held-out checks for each detector model and
  keep out-of-domain behavior explicit.
  \item Integrate the native desktop shell with versioned scientific services,
  run state, persistence, event views, and platform-specific acceptance work.
  \item Complete accessible learning materials, human review, packaging,
  preservation, and a fresh-environment release review.
\end{enumerate}

Detailed requirements, priority assignments, and acceptance evidence are
maintained in \path{docs/engineering/requirements.json},
\path{docs/engineering/implementation-evidence.json}, and the numbered
engineering handbook chapters. A roadmap item should be considered complete
only when its implementation and evidence agree.

\section{Getting started}

The following commands reflect the repository's documented development entry
points. Optional dependencies and external runtimes are needed for specific
workflows.

\subsection{Python analysis package}

\begin{lstlisting}
pip install -e .
pip install -e '.[analysis]'

heplab version
heplab catalog list
heplab run dimuon --input events.json --output artifacts/dimuon.json
heplab plot artifacts/dimuon.json
\end{lstlisting}

The dimuon input may be a JSON list of schema-v1 events or an object containing
an \texttt{events} list; the reader also supports NanoAOD ROOT input. Plotting,
fitting, parquet caching, and ROOT reading may require optional dependencies.
Consult the operations handbook before acquiring data or relying on a
particular runtime.

\subsection{Local interfaces and documentation}

\begin{lstlisting}
heplab gui --open-browser
dotnet run --project web/HEPLab.Web
mkdocs build --strict
cmake --preset windows-ninja
\end{lstlisting}

These commands refer to separate interfaces: the first opens a local
compatibility workbench, the second runs the ASP.NET Core site, the third
builds the engineering handbook, and the fourth configures the native desktop
preset. Successful setup depends on the relevant SDKs and system dependencies.
A configured or built shell does not by itself indicate complete backend
integration or physics validation.

\section{Licensing, data use, and attribution}

The HEPLab source is licensed under Apache-2.0; the repository's
\path{LICENSE} is authoritative. CERN Open Data records have their own
record-level terms and should be cited by their persistent identifiers and
exact data provenance. A result should preserve the record, file or entry
subset, and certification payload used.

External software and simulation content have separate terms. In particular,
PYTHIA licensing and citation obligations must be reviewed for any redistributed
binary or container. The HEPLab source license does not replace dependency,
dataset, or generated-content obligations.

\section{Further reading}

This document is a summary. The following repository materials are the
operational and technical sources for specific questions:

\begin{itemize}
  \item \textbf{Repository overview:} \path{README.md}.
  \item \textbf{Engineering handbook and current delivery plan:}
  \path{docs/engineering/README.md},
  \path{docs/engineering/implementation.md},
  \path{docs/engineering/17-delivery-plan.md}.
  \item \textbf{Scientific assessment and validation:}
  \path{docs/scientific-review/01_PROJECT_OVERVIEW.md},
  \path{docs/scientific-review/02_PHYSICS_VALIDATION.md},
  \path{docs/scientific-review/10_EXECUTIVE_ASSESSMENT.md}.
  \item \textbf{CMS data and analysis:}
  \path{docs/engineering/02-data-supply.md},
  \path{docs/engineering/09-analysis.md},
  \path{docs/engineering/cms-physics-regression-v2.md},
  \path{docs/engineering/19-sources.md}.
  \item \textbf{Simulation and educational models:}
  \path{docs/engineering/05-generation.md},
  \path{docs/engineering/06-transport.md},
  \path{docs/engineering/07-detector-response.md},
  \path{docs/engineering/material-response.md},
  \path{docs/engineering/learning-platform.md}.
  \item \textbf{Desktop status:}
  \path{docs/engineering/desktop-rebuild-audit.md},
  \path{docs/engineering/gui-design.md}.
  \item \textbf{External starting points:}
  \href{https://opendata.cern.ch/record/30522}{CMS Open Data record 30522},
  \href{https://opendata.cern.ch/record/14220}{2016 certified-run record 14220},
  \href{https://pythia.org/documentation/}{PYTHIA documentation},
  \href{https://delphes.github.io/workbook/}{DELPHES workbook}.
\end{itemize}

\end{document}

